% part: intuitionistic-logic
% chapter: propositions-as-types
% section: reduction

\documentclass[../../../include/open-logic-section]{subfiles}

\begin{document}

\olfileid{int}{pty}{red}

\olsection{ઘટાડો}

સ્વાભાવિક નિગમનની !!{derivation}sમાં
!!a{introduction} નિયમ પછી
!!a{elimination} નિયમ આવે, તો તે
આડમાર્ગ છે જેને ``ઘટાડી'' શકાય છે.
દાખલા તરીકે, $\Intro{\land}$ પછી
$\Elim{\land}$ આવે ત્યારે તે
``એકબીજાને રદ કરે છે'': \Intro{\land}થી
બનતા સંયોજનના બે સંયોજ્યો પૈકી
એક $\Elim{\land}$નો નિષ્કર્ષ હોય છે,
અને $\Intro{\land}$નાં આધારવિધાનો એ
બંને સંયોજ્યો છે. તેથી કોઈ પણ
સંયોજ્યની !!a{derivation} જોઈએ, તો
સંયોજન દાખલ કરીને પછી તેને દૂર
કરવાને બદલે સીધી એ જ !!{derivation}
વાપરી શકીએ. આમ:
\begin{gather*}
  \bottomAlignProof
  \AxiomC{}
  \RightLabel{$\delta_1$}
  \DeduceC{$!A_1$}
  \AxiomC{}
  \RightLabel{$\delta_2$}
  \DeduceC{$!A_2$}
  \RightLabel{$\Intro{\land}$}
  \BinaryInfC{$!A_1 \land !A_2$}
  \RightLabel{$\Elim{\land}_i$}
  \UnaryInfC{$!A_i$}
  \DisplayProof\bottomAlignProof
  \quad
  \redone
  \quad
  \AxiomC{}
  \RightLabel{$\delta_i$}
  \DeduceC{$!A_i$}
  \DisplayProof
\end{gather*}

આ સરળીકરણ પરથી સંબંધિત સાબિતી-પદો પર
ઘટાડાનો સમાન સંબંધ લઈ શકીએ.
$\delta_1$નું સાબિતી-પદ~$N_1$ અને
$\delta_2$નું પદ~$N_2$ હોય, તો
ડાબી સાબિતીનું પદ એક પગલામાં
જમણી સાબિતીના પદમાં ઘટે છે:
\[
  \proj{i}{\pair{N_1}{N_2}} \redone N_i
\]
પ્રકારિત લૅમ્બડા-કલનમાં આ
ગુણાકાર-પ્રકારનો \emph{બીટા-ઘટાડાનો}
નિયમ છે.
\sourcecorrection{OLMD-291}{મૂળ જોડીના બંને પદ એક જ આર્ગ્યુમેન્ટમાં લખે છે. સાબિતી-પદની વ્યાખ્યા અને નિયમ-કોષ્ટક મુજબ જોડીના બંને આર્ગ્યુમેન્ટ અલગ કર્યા છે.}

સ્વાભાવિક નિગમનમાં \Intro{\land} પછી
\Elim{\land} જેવા બે અનુમાનો, એટલે
ઘટાડી શકાય એવાં અનુમાનોની જોડી,
\emph{કટ} કહેવાય છે. પ્રકારિત
લૅમ્બડા-કલનમાં $\redone$ની ડાબી
બાજુનું પદ \emph{સંકોચ્ય પદ} અને
જમણી બાજુનું પદ \emph{સંકોચનનું પરિણામ}
કહેવાય છે. અપ્રકારિત લૅમ્બડા-કલનમાં
ફક્ત $(\lambd[x][N])Q$ સંકોચ્ય પદ
ગણાય છે; પ્રકારિત લૅમ્બડા-કલનમાં
વાક્યરચનામાં $\pair{N}{M}$,
$\proj{i}{N}$, $\inj[!A]{i}{N}$,
$\dcase{N}{x_1}{M_1}{x_2}{M_2}$ અને
$\abort{!A}{N}$ ધરાવતાં પદો પણ
ઉમેરાય છે, જેમનાં અનુરૂપ સંકોચ્ય પદો છે.
\sourcecorrection{OLMD-292}{મૂળ વિકલ્પ-દાખલના ચિહ્ન માટે ત્રણ ફરજિયાત આર્ગ્યુમેન્ટ લખે છે; વ્યાખ્યા અને નિયમ-કોષ્ટક મુજબ પ્રકારની નોંધ વૈકલ્પિક આર્ગ્યુમેન્ટમાં મૂકી છે.}

એ જ રીતે વિકલ્પ માટે ઘટાડો છે:
\begin{gather*}
  \bottomAlignProof
  \AxiomC{}
  \RightLabel{$\delta$}
  \DeduceC{$!A_i$}
  \RightLabel{$\Intro{\lor}$}
  \UnaryInfC{$!A_1 \lor !A_2$}
  \AxiomC{$\Discharge{!A_1}{x_1}$}
  \RightLabel{$\delta_1$}
  \DeduceC{$!C$}
  \AxiomC{$\Discharge{!A_2}{x_2}$}
  \RightLabel{$\delta_2$}
  \DeduceC{$!C$}
  \DischargeRule{$\Elim{\lor}$}{x_1,x_2}
  \TrinaryInfC{$!C$}
  \DisplayProof\bottomAlignProof
  \quad
  \redone
  \quad
  \AxiomC{}
  \RightLabel{$\delta$}
  \DeduceC{$!A_i$}
  \RightLabel{$\delta_i$}
  \DeduceC{$!C$}
  \DisplayProof
\end{gather*}
તે સાબિતી-પદો પરના આ ઘટાડાને અનુરૂપ છે:
\begin{gather*}
  \dcase{\inj[!A_{3-i}]{i}{M}}{x_1}{N_1}{x_2}{N_2} \redone \Subst{N_i}{M}{x_i}
\end{gather*}
અહીં $M$ એ~$\delta$નું સાબિતી-પદ અને
$N_i$ એ~$\delta_i$નું પદ છે.
આ સરવાળા-પ્રકારનો બીટા-ઘટાડાનો નિયમ છે.
અહીં $\Subst{M}{N_i}{x}$ એટલે~$M$માં
ચલ~$x$ની દરેક મુક્ત આવૃત્તિને~$N_i$થી
બદલવી.
\sourcecorrection{OLMD-293}{મૂળ વિકલ્પ-દાખલની વૈકલ્પિક નોંધમાં દાખલ થતા પદનો પોતાનો પ્રકાર આપે છે. નિયમ-કોષ્ટક પ્રમાણે આ નોંધ બીજો, હજી નિર્દિષ્ટ ન થયેલો વિકલ્પ આપે છે; તેથી અહીં $!A_{3-i}$ કર્યું છે.}

વિધેય-પ્રકારનો ઘટાડો $\Intro\lif$
પછી $\Elim\lif$ના આડમાર્ગને
દૂર કરવાને અનુરૂપ છે.
\begin{gather*}
  \bottomAlignProof
  \AxiomC{$\Discharge{!A}{x}$}
  \RightLabel{$\delta$}
  \DeduceC{$!B$}
  \DischargeRule{$\Intro{\lif}$}{x}
  \UnaryInfC{$!A \lif !B$}
  \AxiomC{}
  \RightLabel{$\delta'$}
  \DeduceC{$!A$}
  \RightLabel{$\Elim{\lif}$}
  \BinaryInfC{$!B$}
  \DisplayProof\bottomAlignProof
  \quad
  \redone
  \quad
  \AxiomC{}
  \RightLabel{$\delta'$}
  \DeduceC{$!A$}
  \RightLabel{$\delta$}
  \DeduceC{$!B$}
  \DisplayProof
\end{gather*}
સાબિતી-પદો માટે આ સામાન્ય બીટા-ઘટાડો છે:
\begin{gather*}
  (\lambd[\typeof{x}{!A}][N]M)
  \redone \Subst{N}{M}{x}
\end{gather*}

!!{proof}sનાં ઘટાડા-રૂપાંતરો ઉપરાંત
ક્રમચય-રૂપાંતરો પણ લઈ શકીએ.
તે એવી (ઉપ-)!!{proof}s પર લાગુ
થાય છે જેના અંતે \Elim{\lor}
પછી બીજો !!{elimination} નિયમ
આવે છે. દાખલા તરીકે:
\begin{multline*}
  \bottomAlignProof
  \AxiomC{}
  \RightLabel{$\delta$}
  \DeduceC{$!A_i$}
  \RightLabel{$\Intro{\lor}$}
  \UnaryInfC{$!A_1 \lor !A_2$}
  \AxiomC{$\Discharge{!A_1}{x_1}$}
  \RightLabel{$\delta_1$}
  \DeduceC{$!C_1 \land !C_2$}
  \AxiomC{$\Discharge{!A_2}{x_2}$}
  \RightLabel{$\delta_2$}
  \DeduceC{$!C_1 \land !C_2$}
  \DischargeRule{$\Elim{\lor}$}{x_1,x_2}
  \TrinaryInfC{$!C_1 \land !C_2$}
  \RightLabel{\Elim{\land}[i]}
  \UnaryInfC{$!C_i$}
  \DisplayProof\bottomAlignProof
  \quad
  \redone
  \\
  \AxiomC{}
  \RightLabel{$\delta$}
  \DeduceC{$!A_i$}
  \RightLabel{$\Intro{\lor}$}
  \UnaryInfC{$!A_1 \lor !A_2$}
  \AxiomC{$\Discharge{!A_1}{x_1}$}
  \RightLabel{$\delta_1$}
  \DeduceC{$!C_1 \land !C_2$}
  \RightLabel{\Elim{\land}[i]}
  \UnaryInfC{$!C_i$}
  \AxiomC{$\Discharge{!A_2}{x_2}$}
  \RightLabel{$\delta_2$}
  \DeduceC{$!C_1 \land !C_2$}
  \RightLabel{\Elim{\land}[i]}
  \UnaryInfC{$!C_i$}
  \DischargeRule{$\Elim{\lor}$}{x_1,x_2}
  \TrinaryInfC{$!C_i$}
  \DisplayProof
\end{multline*}
$\delta$, $\delta_1$ અને $\delta_2$નાં
સાબિતી-પદો અનુક્રમે $M$, $N_1$ અને
$N_2$ હોય, તો સાબિતી-પદો, એટલે
$\lambd$-પદો, માટેનો અનુરૂપ
ઘટાડા-નિયમ છે:
\[
  \proj{i}{\dcase{M}{x_1}{N_1}{x_2}{N_2}} \redone 
\dcase{M}{x_1}{\proj{i}{N_1}}{x_2}{\proj{i}{N_2}}
\]

આડમાર્ગો ફક્ત !!{proof}sના અંતે જ
નહીં, પણ !!{proof}sની અંદર પણ
દૂર કરી શકીએ: આડમાર્ગથી પૂરી થતી
ઉપ-!!{proof} પર ઘટાડા-રૂપાંતર લગાવો.
એ જ રીતે $\beta$-ઘટાડો
$\lambd$-પદોના ઉપપદો પર લાગુ
થાય છે. $N$ એ~$M$નું ઉપપદ હોય
અને $N \redone N'$ થાય, તો
$M$માંનું ઉપપદ~$N$ એ~$N'$થી
બદલી $M'$ મેળવી શકીએ.
ત્યારે પણ $M \redone M'$ લખીએ.
$M = M_1 \redone M_2 \redone M_3 \redone \dots \redone M_k = M'$
અનુક્રમ હોય ($k=1$, એટલે $M = M'$
હોય તે કિસ્સો સહિત), તો
$M \red M'$ લખીએ.
\sourcecorrection{OLMD-294}{મૂળ અનુક્રમમાં બીજું પદ બે વખત લખાયું છે; ક્રમશઃ ત્રીજું પદ મૂક્યું છે.}

જેના કોઈ પણ ઉપ-!!{proof} પર રૂપાંતર
લાગુ ન થઈ શકે તેવી !!{proof}
\emph{સામાન્ય} કહેવાય છે.
એ જ રીતે જેના કોઈ પણ ઉપપદ પર
રૂપાંતર લાગુ ન થઈ શકે એવું
$\lambd$-પદ \emph{સામાન્ય} કહેવાય છે.
સામાન્ય $\lambd$-પદને
\emph{સામાન્ય સ્વરૂપ} પણ કહે છે.

\end{document}
